\subsection{Free cash flow valuation}

\begin{qfmethod}[Free cash flow to the firm and to equity: definitions, calculation from net income, EBIT, EBITDA and operating cash flow, and the ownership perspective each represents]\label{meth:fcff-fcfe-derivation}
\qffield{Statement} \emph{Free cash flow to equity} is the cash left after operating expenses, taxes, debt payments and reinvestment: the cash a firm could pay to its shareholders \citep[pp.~351--352]{investment-valuation-2012}. \emph{Free cash flow to the firm} is the cash flow to all claimholders, before any payment to lenders or preferred shareholders (pp.~380--381):
\[ \mathrm{FCFF}=\mathrm{EBIT}(1-t)+\mathrm{Dep}-\mathrm{FCInv}-\mathrm{WCInv},\qquad \mathrm{FCFE}=\mathrm{FCFF}-\mathrm{Int}(1-t)+\text{net borrowing}. \]
Interest's tax saving is left out of FCFF because the after-tax cost of debt in the WACC already counts it (p.~381). Equivalent routes (Derivation, from the identities $\mathrm{EBIT}=\mathrm{EBITDA}-\mathrm{Dep}$ and $\mathrm{NI}=(\mathrm{EBIT}-\mathrm{Int})(1-t)$):
\begin{align*}
\mathrm{FCFF}&=\mathrm{NI}+\mathrm{NCC}+\mathrm{Int}(1-t)-\mathrm{FCInv}-\mathrm{WCInv}\\
&=\mathrm{EBITDA}(1-t)+\mathrm{Dep}\times t-\mathrm{FCInv}-\mathrm{WCInv}\\
&=\mathrm{CFO}+\mathrm{Int}(1-t)-\mathrm{FCInv},
\end{align*}
where NCC are non-cash charges and CFO deducts interest paid, as US GAAP requires.

\emph{Where interest and its tax saving are classified.} IAS~7 before IFRS~18 let a non-financial entity classify interest paid in operating or financing \citep[IAS~7 para.~33]{eu-2023-1803-ias7}. Once IFRS~18 applies, interest paid is financing for entities whose main business is not investing or lending \citep[IAS~7 para.~34A as amended]{eu-2026-338-ifrs18}. Under both versions, income taxes paid stay in operating unless specifically identified with financing or investing (para.~35). CFO is then before interest but after the actual tax paid, which is already reduced by the interest deduction. The tax saving has to come out (Derivation, from the identities above):
\[ \mathrm{FCFF}=\mathrm{CFO}_{\text{interest in financing}}-\mathrm{Int}\times t-\mathrm{FCInv}. \]
Only if the tax effect of interest were also classified in financing would CFO equal operating cash flow after tax on EBIT, and FCFF be $\mathrm{CFO}-\mathrm{FCInv}$. FCFE equals net income plus non-cash charges, less investment, plus net borrowing (p.~352). Preferred dividends are deducted, and new preferred issued is added, when preferred stock exists (p.~353).

\emph{Ownership perspective.} FCFE measures what equity holders could receive. Dividends are what a minority holder actually receives; FCFE is what a holder who controlled payout could take out. Valuing FCFE when a firm pays less than it could therefore values control of payout policy (pp.~374--375).

\qffield{Notation} $t$ tax rate; Dep depreciation and amortisation; FCInv investment in fixed capital; WCInv investment in working capital; Int interest expense.

\qffield{Practice} \emph{Avon Fittings} (original, one year). Revenue 500, cash costs 380, depreciation 30, interest 12, tax 25 percent, fixed-capital investment 45, working-capital investment 8, net borrowing 10:
\begin{itemize}
\item EBITDA 120, EBIT 90, net income $(90-12)\times0.75=58.50$, CFO $58.5+30-8=80.50$;
\item FCFF from net income: $58.5+30+9-45-8=44.50$; from EBIT: $67.5+30-53=44.50$; from EBITDA: $90+7.5-53=44.50$; from CFO: $80.5+9-45=44.50$;
\item FCFE: $44.5-9+10=45.50$, equal to $58.5+30-53+10$.
\end{itemize}
Taxes paid are 19.50, on profit before tax of 78. With interest paid in financing and taxes in operating, CFO is $120-19.5-8=92.50$, and FCFF is $92.5-12\times0.25-45=44.50$. Deducting only capital investment would give 47.50, crediting the interest tax saving to FCFF a second time. CFO would be 89.50, and FCFF $89.5-45$, only if the tax saving of 3 were also shown in financing. All three presentations describe the same 44.50.

\emph{Harbour, 2026 forecast} (\hyperref[meth:free-cash-flow-coverage]{FCFF and FCFE}). Valuation FCFF of 31.93 follows two routes. One starts from CFO, adds after-tax interest and pension financing, and deducts capital spending, right-of-use additions and share-based pay. The other starts from EBIT less cash operating taxes, adds depreciation and deducts working capital and all reinvestment. Both give the same 31.93.

\qffield{Conditions} The routes agree only when each counts the same reinvestment and the same claims. Leases treated as debt make right-of-use additions reinvestment; share-based pay treated as a cost is not added back. The CFO routes assume taxes paid are classified in operating activities, as IAS~7 para.~35 requires unless they can be identified with another activity, and that interest is fully deductible at the rate $t$ in the year it is paid.

\qffield{Pitfalls} Adding back depreciation and also deducting gross rather than net capital spending inconsistently. Subtracting after-tax interest from FCFF and discounting at the WACC, which counts the tax saving twice. Using CFO without checking where interest, dividends and the related taxes paid are classified; with interest in financing, CFO less capital investment overstates FCFF by the interest tax saving.

\qffield{See also} \hyperref[meth:fcf-forecasting]{forecasting}; \hyperref[meth:cash-flow-interpretation]{cash flow statement}.

\end{qfmethod}

\begin{qfmethod}[Forecasting FCFF and FCFE: forecasting the components or applying a target debt ratio, and what dividends, buybacks, share issues and changes in leverage do to each]\label{meth:fcf-forecasting}
\qffield{Statement} Free cash flows are forecast either by forecasting each component (net income or EBIT, depreciation, capital spending, working capital and borrowing) or, for FCFE, by assuming that a fixed share $\delta$ of net reinvestment is financed by debt \citep[pp.~352--353]{investment-valuation-2012}:
\[ \mathrm{FCFE}=\mathrm{NI}-(1-\delta)(\mathrm{FCInv}-\mathrm{Dep})-(1-\delta)\,\mathrm{WCInv}. \]
The second form keeps leverage stable by construction and smooths the forecast (p.~354).

\emph{Effects of financing decisions.} Neither measure includes dividends, buybacks or share issues, which are uses or sources of equity cash, not operating flows. Because FCFE already contains the net borrowing actually done, its uses reconcile exactly (Derivation, from the cash-flow statement):
\[ \mathrm{FCFE}=\text{dividends}+\text{buybacks}-\text{share issues}+\Delta\text{cash}. \]
Dividends above this FCFE are therefore met by issuing shares or running down cash, not by further borrowing, as Damodaran says; a firm paying less accumulates cash (p.~355). Measured against FCFE at a target debt policy instead, borrowing beyond the policy is a third source of the excess (Interpretation). Borrowing more raises FCFE in the year by the amount borrowed less the after-tax interest, and lowers later FCFE by the interest and repayments. FCFF is unaffected, but leverage raises the cost of equity, so more FCFE does not mean more equity value (p.~361).

\emph{Fixed against endogenous borrowing.} Whether payout changes FCFE depends on the financing policy in the forecast (Interpretation):
\begin{itemize}
\item with net borrowing fixed in advance, a higher dividend reduces cash and leaves FCFE unchanged;
\item with a revolver that tops cash up to a minimum, the higher dividend is borrowed, so actual FCFE rises by roughly the extra dividend and later FCFE falls by the interest.
\end{itemize}
A valuation should therefore use FCFE at a target borrowing policy, which does not move with the payout chosen, and keep actual borrowing for the cash reconciliation.

\qffield{Practice} \emph{Avon Fittings} (\hyperref[meth:fcff-fcfe-derivation]{derivation}). With 30 percent of net reinvestment financed by debt, $\mathrm{FCFE}=58.50-0.70\times(45-30+8)=42.40$. Net borrowing is fixed at 10, so a dividend or buyback changes cash, not FCFF or FCFE: dividends of 30 and a buyback of 10 leave a cash increase of 5.50, and five more of dividends cut it to 0.50. Borrowing 20 more at 6 percent, drawn at the start of the year, raises FCFE by $20-0.90=19.10$ in the year and leaves FCFF at 44.50.

\emph{Harbour, 2026--2028} (\hyperref[ex:harbour-pro-forma-model]{forecast}). The model holds cash at about 40 by drawing on the revolver, and issues and buys back no shares. Actual FCFE therefore equals dividends plus the change in cash: 31.82, 30.19 and 32.14. In 2026, 29.36 of dividends and 2.46 of cash increase use it all. Net borrowing, including new lease debt, is part of FCFE.
\begin{itemize}
\item At a 50 percent payout instead of 40, the 2026 dividend is 36.62 and the revolver draw 38.28 instead of 30.81. Actual FCFE becomes 39.08, again dividends plus the unchanged 2.46 of cash.
\item FCFE before the revolver, with only scheduled loan and lease flows, is 1.01 at 40 percent and 0.80 at 50 percent: the payout moves it only through the extra interest.
\item Valuation FCFF moves from 31.93 to 31.88. FCFF credits the whole tax saving on the extra interest, but the model settles only a fifth of current tax within the year, so most of that saving is not yet received in cash.
\item A share issue would add to cash or replace revolver drawing, so it would change actual FCFE only through borrowing.
\end{itemize}

\qffield{Conditions} The target-ratio shortcut is consistent only if the debt ratio is held in book terms and the cost of equity reflects that ratio. Net borrowing in the forecast must follow a stated policy, not be whatever balances the model. With a revolver that balances the model, actual FCFE measures what the policy paid out, not what the firm could pay.

\qffield{Pitfalls} Treating equity raised as part of FCFE. Forecasting FCFE with large one-off borrowing that is not repeated, then capitalising it. Letting forecast leverage drift while the cost of equity stays fixed. Valuing an FCFE that includes revolver drawings that rise with the dividend, which makes the value depend on the payout chosen. Saying a dividend above actual FCFE was funded by borrowing, when the borrowing is already inside FCFE.

\qffield{See also} \hyperref[meth:fcf-valuation-models]{valuation models}; \hyperref[meth:dividend-sustainability]{dividend sustainability}.

\end{qfmethod}

\begin{qfmethod}[Valuing equity with free cash flow: choosing FCFF or FCFE, single-stage and multistage models, terminal values consistent with stable growth, and sensitivity analysis]\label{meth:fcf-valuation-models}
\qffield{Statement} \emph{Models.}
\begin{itemize}
\item FCFE discounted at the cost of equity values equity directly.
\item FCFF discounted at the WACC values the operating assets; debt and other claims are then deducted and non-operating assets added \citep[pp.~358--362, 383]{investment-valuation-2012}.
\item In stable growth, $V_E=\mathrm{FCFE}_1/(k_e-g)$ and $V_F=\mathrm{FCFF}_1/(\mathrm{WACC}-g)$.
\item A two-stage model adds the present value of high-growth cash flows to a terminal value $P_n=\mathrm{FCFE}_{n+1}/(k_{e,st}-g_n)$ (p.~362).
\item A three-stage model inserts a transition in which growth, the reinvestment rate and the cost of equity move from their high-growth to their stable values, discounting at the cumulated cost of equity (pp.~374--375).
\end{itemize}

\emph{Choosing the number of stages} (Interpretation, consistent with Damodaran's examples): a firm already growing at about the economy's rate needs one stage; a firm with a clear advantage that ends abruptly, such as a patent, two; a firm whose growth will fade gradually as competition arrives, three.

\emph{Choosing between FCFF and FCFE.} FCFE suits firms with stable leverage. FCFF is better when leverage is high or changing, or FCFE is negative, because the WACC is less sensitive to leverage than the cost of equity (Interpretation, consistent with pp.~361 and~402). Both give the same equity value if leverage is held constant and the discount rates are consistent.

\emph{Terminal value.} Stable growth cannot exceed the growth of the economy. The firm must also look stable: capital spending close to depreciation and a beta near one, or reinvestment set at $g/\mathrm{ROE}$ for FCFE and $g/\mathrm{ROC}$ for FCFF (pp.~359--363, 383). Damodaran's example shows that letting capital spending of twice depreciation grow into perpetuity understates the terminal value (pp.~362--363). The alternative, a multiple of terminal-year earnings or EBITDA, imports the market's pricing of peers and must imply a growth and return consistent with the forecast (\hyperref[meth:comparables-evaluation]{terminal values from multiples}).

\qffield{Practice} \emph{Brent Utilities} (original, stable growth with debt held at 40 percent of value). FCFF next year is 60, growth 3 percent, WACC 7 percent, pre-tax cost of debt 5 percent, tax 25 percent:
\begin{itemize}
\item firm value $60/(7\%-3\%)=1{,}500$; debt 600; equity 900;
\item cost of equity $k_e=(0.07\times1{,}500-0.0375\times600)/900=9.17$ percent;
\item FCFE $=60-0.0375\times600+0.03\times600=55.50$, since debt grows with the firm; $55.50/(9.17\%-3\%)=900$, the same equity value, confirmed by an explicit 3,000-year sum;
\item discounting FCFE at the WACC instead would give 1,387.50, a mismatch of cash flow and rate.
\end{itemize}

\emph{Two-stage FCFE} (original, per share). FCFE of 2.00 grows 12 percent for five years, to 3.52 in year 5. The cost of equity is 10 percent in high growth, then 9 percent with 3 percent growth. The terminal value is $3.52\times1.03/6\%=60.51$, and the value is 48.13. Stable growth with ROE of 12 percent needs equity reinvestment of $3\%/12\%=25$ percent.

\emph{Three-stage} (same firm). Growth falls to 9, 6 and 3 percent over years 6--8 while the cost of equity steps down to 9 percent. The terminal value at year 8 is 72.01 and the value 51.05, 2.92 more than the two-stage value: the transition years carry growth that the two-stage model drops at once.

The stable-stage inputs dominate the value:
\begin{center}\small
\begin{tabular}{@{}lrrr@{}}
\toprule
Stable $k_e$ & $g=2\%$ & $g=3\%$ & $g=4\%$ \\
\midrule
8\% & 47.76 & 55.64 & 67.46 \\
9\% & 42.45 & 48.13 & 56.08 \\
10\% & 38.46 & 42.76 & 48.49 \\
\bottomrule
\end{tabular}
\end{center}

\emph{Harbour} uses FCFF because its revolver drawing and lease debt change from year to year (\hyperref[meth:forecast-model-valuation]{forecast model}).

\qffield{Conditions} The constant-leverage reconciliation assumes debt is rebalanced to a fixed share of value, so the interest tax saving carries the firm's operating risk. A firm with a fixed debt amount needs adjusted present value or a year-by-year cost of capital.

\qffield{Pitfalls} Discounting FCFE at the WACC or FCFF at the cost of equity. A terminal growth rate above the economy's. Stable growth with high-growth reinvestment or beta. Reporting a value without its sensitivity to the stable-stage inputs.

\qffield{See also} \hyperref[meth:fcff-fcfe-derivation]{FCFF and FCFE}; \hyperref[meth:multistage-ddm]{multistage dividend models}; \hyperref[meth:forecast-horizon-steady-state]{steady state}.

\end{qfmethod}

\begin{qfmethod}[Cash-flow proxies compared: FCFE against dividend discount models, and net income and EBITDA as stand-ins for free cash flow]\label{meth:cash-flow-proxies}
\qffield{Statement} \emph{FCFE against dividends.} The two models give the same value when dividends equal FCFE, or when the excess cash is invested at a net present value of zero \citep[p.~372]{investment-valuation-2012}. FCFE gives the higher value when excess cash is wasted or held earning less than its cost, and the lower value when dividends above FCFE are financed by costly issues or excess debt (pp.~372--373). The FCFE value is the better estimate when control of payout can change, as in a takeover, and the dividend value when it cannot (p.~374).

\emph{Net income and EBITDA as proxies.} Damodaran relates each to FCFF (p.~382, Table~15.1):
\begin{align*}
\mathrm{EBIT}(1-t)&=\mathrm{FCFF}+\text{capex}-\text{depreciation}+\Delta\mathrm{WC},\\
\mathrm{EBITDA}&=\mathrm{FCFF}+\mathrm{EBIT}\times t+\text{capex}+\Delta\mathrm{WC}.
\end{align*}
\begin{itemize}
\item Discounting after-tax operating income assumes no net reinvestment, so it allows no growth (p.~382).
\item Discounting EBITDA assumes no taxes and that the firm disinvests; an infinite life or growth is then inconsistent (p.~382).
\item Net income ignores reinvestment and borrowing, so it equals FCFE only if net investment is entirely debt-financed.
\end{itemize}

\qffield{Practice} \emph{Harbour, 2026 forecast.}
\begin{itemize}
\item EBITDA 227.28;
\item NOPAT 87.28 (EBIT 119.52 less tax on EBIT plus non-deductible share-based pay);
\item net income 73.40;
\item valuation FCFF 31.93;
\item FCFE 31.82;
\item dividends 29.36.
\end{itemize}
The gap between EBITDA and FCFF is taxes, capital spending of 120.31, right-of-use additions of 25 and share-based pay. Capitalising EBITDA as if it were FCFF, $227.28/(8\%-2.5\%)$, would give 4,132 against a DCF enterprise value of 1,218. Dividends are close to actual FCFE by construction: the revolver borrows whatever the 40 percent payout and the cash target need, so actual FCFE follows the payout (39.08 at 50 percent; \hyperref[meth:fcf-forecasting]{forecasting}). Comparing a dividend model with an FCFE model needs FCFE at a target debt policy, not the actual borrowing.

\emph{Avon Fittings} (\hyperref[meth:fcff-fcfe-derivation]{derivation}): $\mathrm{EBIT}(1-t)=67.50=44.50+45-30+8$ and $\mathrm{EBITDA}=120=44.50+22.50+45+8$.

\qffield{Conditions} The comparisons use the valuation conventions of the Harbour case. For banks and insurers, debt is raw material, and dividend or excess-return models are used instead (\hyperref[meth:bank-camels-analysis]{banks}).

\qffield{Pitfalls} Calling EBITDA cash flow. Valuing a firm by its dividends when it pays far less than its FCFE and has no investor able to change that, or by FCFE when no one can. Using net income growth as FCFE growth when the firm is changing its leverage.

\qffield{See also} \hyperref[meth:enterprise-value-multiples]{EV multiples}; \hyperref[meth:fcf-valuation-models]{valuation models}; \hyperref[meth:dividend-sustainability]{dividend cover}.

\end{qfmethod}
